\documentclass[11pt,a4paper]{article}
\usepackage{fontspec,amsmath,amssymb,amsthm,microtype}
\setmainfont{lmroman10-regular.otf}[BoldFont=lmroman10-bold.otf,ItalicFont=lmroman10-italic.otf,BoldItalicFont=lmroman10-bolditalic.otf]
\setmonofont{lmmono10-regular.otf}[Scale=MatchLowercase]
\usepackage[margin=25mm]{geometry}
\usepackage[unicode,hidelinks]{hyperref}
\setlength{\emergencystretch}{3em}
\setlength{\parskip}{0.4em}
\setlength{\parindent}{0pt}
\allowdisplaybreaks[2]
\raggedbottom
\title{\texorpdfstring{Regroup genuinely independent source blocks}{Regroup genuinely independent source blocks}}
\author{LangevinLean mathematical exposition}
\date{}
\begin{document}\maketitle
\noindent Source module: \texttt{Proximal\allowbreak{}Langevin/\allowbreak{}Independence\allowbreak{}Helpers.\allowbreak{}lean}. The accompanying HTML pairs this mathematical exposition with the exact Lean statements and records the scope of the formalization.


These measure-theoretic lemmas reorganize supplied joint independence. They do not infer joint independence from pairwise independence.


\section{\texorpdfstring{Two independent variables form a Boolean family}{Two independent variables form a Boolean family}}


Given \(X\perp Y\), define \(Z_{\rm false}=X\), \(Z_{\rm true}=Y\). The finite-family intersection criterion has four possible index subsets: the empty set, either singleton, or both indices. Probability mass one handles the empty case, the singletons are identities, and the two-element case is exactly \(P(X\in A,Y\in B)=P(X\in A)P(Y\in B)\).




An independent family is defined by factorization for every finite set of indices. For a Boolean index there are four possibilities: neither variable, only \(X\), only \(Y\), or both. The singleton identities are tautological; the empty family has probability one because \(\mu\) is a probability; the two-variable case is exactly the assumed independence of \(X\) and \(Y\). The two membership case splits in Lean exhaust these four possibilities explicitly.




\textbf{Correspondence with the displayed declarations.}


\begin{itemize}

\item \texttt{i\allowbreak{}Indep\allowbreak{}Fun\_\allowbreak{}bool\_\allowbreak{}of\_\allowbreak{}indep\allowbreak{}Fun}: Two independent random variables form an independent family indexed by \texttt{Bool}. The probability assumption supplies the empty-family identity.
\end{itemize}


{\raggedright\noindent Lean declarations: \texttt{i\allowbreak{}Indep\allowbreak{}Fun\_\allowbreak{}bool\_\allowbreak{}of\_\allowbreak{}indep\allowbreak{}Fun}.\par}

\section{\texorpdfstring{Regroup a jointly independent triple}{Regroup a jointly independent triple}}


Assume \(X\perp(Y,Z)\) and \(Y\perp Z\), with measurable variables. The joint pushforward is \(\mu_X\otimes(\mu_Y\otimes\mu_Z)\). Use the measurable product-associativity equivalence to rewrite it as \((\mu_X\otimes\mu_Y)\otimes\mu_Z\). Projection also proves \(\mu_{(X,Y)}=\mu_X\otimes\mu_Y\), so \((X,Y)\perp Z\). The joint independence premise is essential.




The hypotheses say that \(X\) is independent of the whole pair \((Y,Z)\), and that \(Y\) and \(Z\) are independent. Therefore the triple law is \(\mu_X\otimes(\mu_Y\otimes\mu_Z)\). Product associativity rewrites it as \((\mu_X\otimes\mu_Y)\otimes\mu_Z\), proving independence of \((X,Y)\) and \(Z\). The proof uses equality of pushforward measures under the measurable associator, so this is joint independence; pairwise independence of three unrelated variables would not suffice.




\textbf{Correspondence with the displayed declarations.}


\begin{itemize}

\item \texttt{indep\allowbreak{}Fun\_\allowbreak{}prod\allowbreak{}Mk\_\allowbreak{}left\_\allowbreak{}of\_\allowbreak{}indep\allowbreak{}Fun}: Regroup a jointly independent triple from \texttt{X | (Y,Z)} into \texttt{(X,Y) | Z}. Independence of \texttt{Y} and \texttt{Z} is essential; pairwise independence alone is insufficient.
\end{itemize}


{\raggedright\noindent Lean declarations: \texttt{indep\allowbreak{}Fun\_\allowbreak{}prod\allowbreak{}Mk\_\allowbreak{}left\_\allowbreak{}of\_\allowbreak{}indep\allowbreak{}Fun}.\par}

\end{document}
